\documentclass[11pt,a4paper]{article}
\usepackage[margin=1.8cm]{geometry}
\usepackage{amsmath,amssymb,array,tabularx,fontspec,cancel}
\setmainfont{FreeSerif}
\setlength{\parindent}{0pt}\setlength{\parskip}{0pt}
\renewcommand{\arraystretch}{1.15}
\newcounter{sq}
\newcommand{\sq}[1]{\stepcounter{sq}\par\smallskip\noindent\hangindent=11em\hangafter=1
  \makebox[11em][l]{\textbf{\Roman{sq}.}}#1\par}
\newcommand{\dtab}[2]{\par\smallskip\noindent\hspace*{11em}\begin{tabular}{#1}\hline #2\\ \hline\end{tabular}\par\smallskip\hspace*{11em}}
\begin{document}
\begin{center}
  {\Huge\bfseries The Hope Science Academy}\\[3pt]
  {\small Excellence is our standard not goal}\\[2pt]
  {\small 247-E-1, (Opposite LDA School) Johar Town, Lahore. PH: 0300-8194789, 0322-4157001}
\end{center}
\medskip\noindent
\begin{tabularx}{\linewidth}{|X|X|X|X|}
\hline
\textbf{Student Name:} & \textbf{Roll Num:} & \textbf{Class:} 11th & \textbf{Paper Code:} \\ \hline
\textbf{Subject Name:} Business Maths & \textbf{Time Allowed:} & \textbf{Total Marks:} 818 & \textbf{Exam Date:} 02-Nov-2026 \\ \hline
\multicolumn{4}{|l|}{\textbf{Exam Syllabus:} FULL BOOK} \\ \hline
\end{tabularx}
\bigskip
\noindent\textbf{Q1. Write short answers of the following questions.}\hfill\textbf{1$\times$818=818}

\sq{What basic arithmetic operations.}
\sq{Define ratio.}
\sq{Define rate.}
\sq{Write two uses of ratio?}
\sq{Differentiate ration and rate.}
\sq{Find the ratio between 5 minutes and 30 seconds.}
\sq{Find the ration between one hour and 30 seconds.}
\sq{Express in reduce from 1.3 to 5.2.}
\sq{Find ration between Rs. 145 and Rs. 105.}
\sq{Simplify the ratio. 4:12:16.}
\sq{Write the ratio $\dfrac{1}{10}$:100 in simplified form.}
\sq{Reduce to its lowest term. 60:180.}
\sq{Express 35:20:45 in lowest form.}
\sq{Find the ratio between one hour 15 minutes and 30 minutes.}
\sq{Find the ratio of 150 centimeter to 3 meters.}
\sq{Simplify the ratio $\dfrac{0.5}{0.06}:\dfrac{7}{8}$ in reduced form.}
\sq{Express in reduce form: 24:64.}
\sq{What is the ratio between 4.5 kg and 7.5 kg.}
\sq{The cash price of Rs, 5000 has to be distributed between two student in the ratio 2:3. How much does each student receive?}
\sq{Divide Rs.60000 in the ratio 5:7.}
\sq{Divide Rs. 3000 in Usman and Bilal in the ratio of 4:6.}
\sq{Distribute Rs. 15000 in the ratio 3:2.}
\sq{Divide 750 in ratio 3:2.}
\sq{Distribute a stock of 6000 electric fans to there dealers in ratio 3:5:4.}
\sq{Divide Rs. 4800 in the ratio 5:7.}
\sq{Alocate Rs. 15000 in ratio 4:6.}
\sq{Divide Rs. 3600 between A and B in the ratio of 3:2.}
\sq{Express 48.5\% as a common fraction.}
\sq{What is the ratio between 128 kg and 16 kg.}
\sq{Three man invested Rs. 18000m Rs.12000 and Rs. 6000 respectively. How should they share out of profit of Rs. 3600?}
\sq{The saving income ratio of two persons are $\dfrac{1}{9}$ and $\dfrac{2}{3}$ who is saving more?}
\sq{Reduce 40 : 240 to its simplest form.}
\sq{Express $\dfrac{4}{9}$ to $\dfrac{1}{3}$ in its simplest form.}
\sq{The cash price of Rs. 1000 has to be distributed between the two students in the ratio 2:3. How much shall each student receive?}
\sq{Bata shoe company has 50,000 pairs of shoes and management like supply these pairs of shoes to four wholesale dealers in the ratio 4:7:9:5. How much pairs of shoes should each dealr receive?}
\sq{If it is given that manager of a particular firm is earning Rs. 100,000 p.m as salary and his spending are Rs. 90000 p.m. One of the subordinates of the same manager is earning Rs. 3000 p.m and he is saving Rs. 6000 p.m. Who is saving more as regard to their income?}
\sq{Define proportion.}
\sq{Find the value of x from 35:125:7:x.}
\sq{What are the different types of proportions?}
\sq{Define direct proportion.}
\sq{Define inverse proportion.}
\sq{Define continued ration.}
\sq{Define compound proportion.}
\sq{What is the basic rule of proportion.}
\sq{Given that 4:x::5:20, find the value of x.}
\sq{Find the value of x if x:250::4:50.}
\sq{Find x if x:$\dfrac{1}{4}$::12:3.}
\sq{Find x if 45:60::900:x.}
\sq{Find the term in the equation of x and y, 2:y::3:x.}
\sq{If 39:x::18:6 then x=?}
\sq{Find fourth proportion of 5,15 and 2.}
\sq{Find the value of x in the proportion 9:x=24:8.}
\sq{Find fourth proportion to 24:18:8.}
\sq{Find x if 14:19::x:38 using its equation form?}
\sq{Find x if 15:60=900:x.}
\sq{Find x if x:$\dfrac{1}{2}$::12:3.}
\sq{Find fourth proportion of number 2,5 and 8.}
\sq{If x:4 :: 9:12. Find the value of x.}
\sq{Find the value of x if 3:4::x:12.}
\sq{Find the missing quantity $\dfrac{9}{x}=\dfrac{30}{270}$}
\sq{Given that x:5::7:2. find the value of x.}
\sq{Find the missing quantity x:50::2:4.}
\sq{Find a:b:c if a=13:5 and b:c = 11:12.}
\sq{Find a:b:c if a:b =17:32 and b:c=48:51.}
\sq{If 15 dozens of eggs cost Rs. 750, how much will 8 dozen of eggs cost?}
\sq{8 pencils cost is Rs. 26. find the cost of 12 pencils.}
\sq{If the price of 50 shirts Rs. 3652. What will be the price of 80 such shirts?}
\sq{If 7 workers paint a house in 6 days. How many days are required to paint the same house if only 3 workers are available?}
\sq{If 20 man can prepare 10 office tables in a day. How many men are required to prepare 25 such office tables in a day?}
\sq{If 15 dozens of eggs cost Rs.750. How much will 8 dozens of eggs cost?}
\sq{If 20 men can construct a housing unit in 60 days. How many men are required to construct such housing unit in 48 days?}
\sq{Rawal express takes four and half hours at a speed of 80 km per hour to go from Rawalpindi to Lahore. How fast the train should travel to complete the same journey in: (i) 3 and half hours (ii) 5 hours}
\sq{Fifteen men are capable to complete a job in 20 days. How much is required to complete the job by 10 men?}
\sq{A production manager plans to produce 100 units with the help of 25 workers, who work 4 hours a day. How many unites 40 workers can make if they work 3 hours per day?}
\sq{If the stay of 12 persons for 28 days in a hotel cost Rs.67200, find the cost for stay of 8 persons for 14 days in the same hotel.}
\sq{Mr, khan , Mr Furqan and Mr. Touseed are three partners they earned a profit of Rs. 18000. The profit can be shared in the ratio: A:B =2:5 and B:C=10:4 Find the share of each partner in profit.}
\sq{Define percentage?}
\sq{Express 20\% as a fraction.}
\sq{Express the 14\% as common fraction.}
\sq{Convert 1.12 into \%.}
\sq{Express $\dfrac{17}{50}$ as a percent.}
\sq{Express 2\% as a fraction.}
\sq{Convert in common fraction 65\%.}
\sq{Express the decimal 0.17 as percent.}
\sq{Express 0.325 as percent.}
\sq{1Express 48.5\% to common fraction.}
\sq{Express common fraction $\dfrac{3}{5}$ as percentage.}
\sq{Express 25\% as a fraction.}
\sq{Express the percentage into decimal $2\dfrac{1}{2}\%$.}
\sq{Express 0.125 as a percentage.}
\sq{Express 2.5\% as a decimal.}
\sq{Calculate Zakat on an amount of Rs. 500000.}
\sq{75 is 25\% of what amount?}
\sq{What number is 20\% more than 900.}
\sq{250 is 20\% of what?}
\sq{A man spend 96\% of his income and save Rs. 525. what is his income?}
\sq{What is 3\% of 20\%.}
\sq{200 is 10\% of what number?}
\sq{Find 10\% of Rs. 15000.}
\sq{Find 20\% of Rs.325000.}
\sq{Rs 250 is $2\dfrac{1}{2}\%$ of what amount?}
\sq{270 is what percent of 900?}
\sq{360 is 30\% of what amount?}
\sq{What percentage is 287 of 700?}
\sq{What is 65\% of 1300?}
\sq{500 is what percentage of 10000?}
\sq{Find 55\% of 800.}
\sq{Find 30\% of 800.}
\sq{Shahid got 680 marks out of 850 marks. Find percentage of his marks?}
\sq{Find 25\% of 12000.}
\sq{What percentage of Rs. 92 is of Rs. 400.}
\sq{500 is 20\% of what amount?}
\sq{What percent of 90 is 22.5.}
\sq{25\% of an amount is 500. find the amount.}
\sq{Find total amount whose 3\% is 1500.}
\sq{Find 30\% of Rs.6000.}
\sq{Find 60\% of 5000.}
\sq{12 is what percentage of 192?}
\sq{30\% of an amount is Rs. 876. Find that amount.}
\sq{Find out 50\% of Rs 2000?}
\sq{240 is what percent of 48000.}
\sq{180 is 20\% of what amount?}
\sq{600 is 10\% of what amount?}
\sq{What percentage of Rs. 120 is 84?}
\sq{Calculate 45\% of 900.}
\sq{What \% is 320 of 800?}
\sq{Find 25\% of 800 marks.}
\sq{5\% of amount is 200, find the amount?}
\sq{160 is 20\% of what amount?}
\sq{Ali got 900 marks out of 1100 marks. Find his percentage of marks.}
\sq{4.6\% of a number is 460. What is the number.}
\sq{Find 20\% of 4000.}
\sq{670.8 is what percentage of 13000?}
\sq{72 is what percentage of 360.}
\sq{525 is what percentage of 10000?}
\sq{Find 40\% of 300.}
\sq{600 is 40\% of what amount?}
\sq{500 is 20\% of what amount?}
\sq{Find 30\% of 258000.}
\sq{Find $13\dfrac{1}{4}\%$ of Rs.400.}
\sq{Translate the statemen ``a\% of b is c'' into an equation.}
\sq{Express 3.2 as percentage.}
\sq{200 is 10\% of what amount?}
\sq{120\% of what number is 200?}
\sq{250 is what percentage of 10000.}
\sq{Convert 15\% to a decimal.}
\sq{Convert 16.25\% to a decimal.}
\sq{Convert 4.9\% to a decimal.}
\sq{Find 20\% of Rs 1100.}
\sq{Find $2\dfrac{1}{2}\%$ of Rs 6635.}
\sq{250 is 20\% of what?}
\sq{60 is what percentage of 360?}
\sq{Define discount.}
\sq{What are the types of discount?}
\sq{Define trade discount.}
\sq{Define quantity discount.}
\sq{Define discount rate.}
\sq{Write formula of calculating discount rate.}
\sq{Write formula of calculating discounted price.}
\sq{Write formula of calculating stated price if discounted price and discount rate are available.}
\sq{Explain the term successive discount?}
\sq{Write the formula of final discounted value with three discounted DR1, DR2 and DR3.}
\sq{Give the formula of computing discount rate for given quantity discount offer.}
\sq{A T.V set of price Rs. 10000 is available for sale at discount of Rs. 500. Give the purchaser price also the discount rate.}
\sq{A retailer received a discount of 15\% on the book value of steel cabinet and paid Rs. 17000. What was the book value of the cabinet?}
\sq{A photo-state machine operator purchased 20 packets of papers. The market price of a packet is Rs, 250. he got a discount @ 17 1/4 and and additional discount @ 4\% for cash payment. Find the final value of the deal}
\sq{Convert the quantity discount buy two get three into equivalent discount rate.}
\sq{What is commission?}
\sq{property agent sold a plot for 93 lac and received 2\% commission. Find the amount received by the agent.}
\sq{What is commission on Rs. 3000 at rate of $33\dfrac{1}{3}\%$}
\sq{Explain commission rate.}
\sq{How commission is different from wage?}
\sq{Brokerage paid on RS. 3500 was Rs. 175, the rate of brokerage is?}
\sq{At $7\dfrac{1}{2}\%$ rate of commission a sales agent got Rs. 4500 on the sales of steel almera. If the price of a steel almera is Rs. 2500. Find the value of his sales and the number of almera. he sold.}
\sq{A broker received Rs. 16000 for selling a piece of land. His rate of commission was $2\dfrac{1}{2}\%$. what was the selling price?}
\sq{Define profit.}
\sq{Define loss.}
\sq{Define profit and loss.}
\sq{Define mark-up.}
\sq{Define markdown.}
\sq{What is rate of profit.}
\sq{What is rate of loss?}
\sq{What is selling prices.}
\sq{A trader bought an electric fan for Rs, 500 and make a profit 121/2. Find the sale price.}
\sq{Find sale value if purchase value is Rs. 4000 and loss is Rs. 359.}
\sq{If cost a bag is Rs. 120 and selling price is Rs. 150 what is the profit or loss.}
\sq{A chain that cost Rs. 200 is sold for Rs. 250 what is the percentage of profit.}
\sq{Aslam bought a cycle fro rs. 15500 and sold it for Rs. 16740. what will be profit or loss in percentage?}
\sq{A radio was sold for Rs. 400, on 10\% loss. find the cost price of radio?}
\sq{If cost of the chair is Rs. 187.50 and is sold for Rs. 250 find percentage profit.}
\sq{A calculator was sold for Rs. 1000 on 5\% loss. Find xost of the calculator.}
\sq{A T.V set was sold for Rs 4000 on 5\% loss. Find the cost of TV.}
\sq{A mobile phone sold for Rs. 8000, if rate of profit is 5\% of sale price, then calculate the sale price of mobile phone.}
\sq{If cost of an article is Rs. 8000 and sale value is Rs. 11500. Find profit percentage.}
\sq{Find sale price if cost is 4000 and loss is 10\%.}
\sq{A man bought a second hand motor car for Rs. 160000 and sold it for Rs. 165000 how much profit or loss did he make?}
\sq{A property dealer purchased a house in Rs. 565000 and sold it for Rs. 525000 in following year.How much profit or loss did he make?}
\sq{Find the gain and selling price of a T.V set of Rs. 4000. sold at 20\% profit.}
\sq{A chair that cost Rs. 187.50 is sold for Rs. 250. Find percentage of profit.}
\sq{An accidential motor cycle of cost Rs.29,500 is sold at a loss of 43.5\%. Find the loss and selling price.}
\sq{A table cost Rs. 600 is sold for Rs. 555, find the percentage loss.}
\sq{An air conditioner was sold at a profit of Rs. 2200. The rate of profit was $12\dfrac{1}{2}\%$. what was the cost of the air conditioner?}
\sq{A jewellery set was sold for rs. 50400 on 16\% loss. Find the cost of the set.}
\sq{Find the cost price if sale price =Rs. 100, loss=25\%.}
\sq{What is meant by interest.}
\sq{Define simple interest.}
\sq{Define principal.}
\sq{Find the simple interest on Rs. 100000 borrowed for 3 years at 5\%.}
\sq{Write formula for simple interest.}
\sq{Find the simple interest on Rs. 8000 invested for 5 years at 9\% per annum.}
\sq{Find simple interest on Rs. 5000 for 10 years at 8\% per annum.}
\sq{What will be the simple interest on Rs 1000 for 5 years at the rate of 10\% annually?}
\sq{Find the amount of simple interst on Rs. 5000 for 3 years at the rate of 4\$ per annum.}
\sq{Find simple interest on Rs. 3000 for 2 months at the rate of 9\% per annum.}
\sq{Find simple interest on Rs. 500 invested for 6 months at 8 \% per annum.}
\sq{Find simple interest of 700 invested for5 years at 5\% per annum.}
\sq{Find the simple interest on Rs. 7000 invested for 5 years at 11\% per annum.}
\sq{Find out simple interest if Rs.8500 are invested for 5 years at 6\% p.a.}
\sq{Find the simple interest on Rs. 3000 at the rate of 8\% per for 5 years.}
\sq{Find simple interest on Rs.500 borrowed for 4 years at 11\% p.a.}
\sq{What principal is needed so that the interest will be Rs. 48 if it invested at 3\% p.a for 5 years.}
\sq{How long will it take for Rs 50000to earn simple interest as Rs. 10,000 at 10\% per year.}
\sq{Find simple interest on Rs. 60000 borrowed for 3 years at 8\% p.a.}
\sq{Find simple interest on Rs. 100000 borrowed for 4 years at 12\% annually,}
\sq{Find the simple interest on Rs. 100000 borrowed for 4 years at 5\%.}
\sq{Find the simple interest on Rs. 7000 invested for 3 years at the rate 11\% annually.}
\sq{Find simple interest on Rs. 100000 borrowed for 3 years at 5\%.}
\sq{Find the simple interest of Rs. 3000 for 2 years at 6\% p.a.}
\sq{Find the simple interest on Rs 700 invested 3 years at 6\% p.a.}
\sq{Find simple interest on Rs. 5000 @ 6\% of 4 years.}
\sq{Find the simple interest on Rs. 4800 at the rate of 11\% for 6 months.}
\sq{Find the simple interest in Rs. 3000 at the rate of 8\% per year for 5 years.}
\sq{Find simple interest of Rs. 10000 borrowed for 4 years at 10\% p.a.}
\sq{Find the simple interest on Rs. 5000 borrowed for 4 years at 11\% annum.}
\sq{Find simple interest on Rs. 2000 for six months 10\% per annum.}
\sq{Find simple interest on Rs. 5000 invested for 4 years at the rate 9\% annually.}
\sq{Find the simple interest on Rs 2000 for one year at 10\%.}
\sq{Find the principal amount ff the simple interest is Rs 500 for 5 years at the rate of 5\%.}
\sq{Find the principal P if the interest for two years at 5\% is Rs. 4000.}
\sq{At what rate Rs. 500 double itself in 5 years by simple interest.}
\sq{Find the rate if Rs. 60000 borrowed for 5 years and simple interest is 30000.}
\sq{How long will it take to earn Rs. 11250 as simple interest on the deposit of Rs. 75000 at 5\% p.a?}
\sq{Find time for Rs. 2000 to become Rs. 2360 at 6\% interest.}
\sq{How long will it take to earn Rs. 15000 interest on a deposit of Rs. 750000 at 10\%.}
\sq{At what rate of 100000 doubled itself in 5 years.}
\sq{At what rate Rs 250 double itself in 5 years by simple interest.}
\sq{How long will it take Rs, 35000 to produce Rs. 7000 as interest at the rate of 4\%.}
\sq{How long will it take for Rs. 50000 to earn simple interest as Rs. 10000 at 10\% per annum.}
\sq{In how many years, simple interest of Rs. 15000can be earned on Rs. 50000 at 10\% yearly?}
\sq{The interest on a loan of Rs. 840 for 9 months was 63. find the rate.}
\sq{If the simple interest on Rs. 1500 for 3 years is Rs. 900 find the rate of interest.}
\sq{Find the total amount at the end of 5 years if a sum of Rs. 200000 borrowed at 13\% annually in simple interest.}
\sq{A bank has issued 3 years loan of rs. 50000. simple interest is charged at the rate of 10 percent per year. The principal plus interest is to be repaid at the end of the third year. Compute the interest for the 3 years period. What amount will be repaid at the end of the third year.}
\sq{A person lends Rs. 100000 to a corporation by purchasing a bond from the corporation. Simple interest is computed quarterly at a rate of 3 percent per quarter. and a cheque for the interest is mailed each quarter to all bondholders. The bond expire at the end of 5 years. and the final cheque includes the original principal plus interest earned during the last quarter. Compute the interest earned each quarter and the total interest. Which will be earned over the 5 year life of the bonds.}
\sq{Define compound interest.}
\sq{Define compound amount.}
\sq{write atleast two points of compound interest.}
\sq{What is the difference between simple interest and compound interest?}
\sq{Write the formula for (a) compound interest (b) Compound interest.}
\sq{Find the compound amount at the end of one year if rs. 2000 are invested at 10\% interest compounded annually.}
\sq{Find the compound amount at the end of one year if Rs. 10000 are invested at 10\% interest compounded annually.}
\sq{Find the compound interest due in case of Rs. 1000 loaned for 5 years at 6\% annually.}
\sq{Find compound interest on Rs. 500 for 6 years at the rate 20\% annually?}
\sq{Calculate the compound interest when Rs 2500 invested for 15 years at 11\% per annum.}
\sq{Find compound amount at the end of 6 years. If Rs.20000 loaned at 8\% compounded annually.}
\sq{Find compound amount when Rs 750 is invested for 8 years at 12\% p.s.}
\sq{Calculate compound interest earned on Rs. 5000 invested for 6 years at the rate of 7\% per annum.}
\sq{At what rate Rs. 5000 double itself in 5 years?}
\sq{Rs. 3000 amounts to Rs. 5843.70 in 17 years compounded annually, what is the interest rate?}
\sq{Define annuity.}
\sq{Define ordinary annuity.}
\sq{Write formula for the amount of an annuity due.}
\sq{Write any two types of annuity.}
\sq{What is amount or maturity value.}
\sq{What are the characteristics of an annuity.}
\sq{Define annuity due.}
\sq{Write formula for the amount of an ordinary annuity.}
\sq{Differentiate between ordinary annuity and annuity due.}
\sq{Write the name of any two field in which annuity used.}
\sq{Define perpetuity.}
\sq{What is amount of an annuity of Rs. 500 payable atan end of each year for 5 years at rate 5\% per annum.}
\sq{An annuity of Rs. 5000 payable at the end of each year amount to Rs. 27082 in 5 years, Find rate?}
\sq{Find the amount of an annuity of Rs 1200 payable at the end of each year for 15 years at the rate of 8\% compounded annually.}
\sq{Find the amount of an annuity of Rs. 1200 payable at the beginning of each year for 15 years @ 8\% compounded annually.}
\sq{How many annual payment of Rs 1000 each be made in the form of an ordinary annuity amounting to Rs. 13180.80 @ 6\% p.a?}
\sq{A father of a girl planned to save some amount monthly to meet the marriage expenditure of his girl. Who is now 15 years old. The father foresight marriage of his girl at the age of 21 years with an amount of Rs 200000 if the money grow at the rate of 12\% p.a compounded monthly then find how many amount he has to save each month.}
\sq{An annuity of Rs. 5000 payable at the end of each year amounts to Rs 27082 in 5 years. find the rate of interest of the annuity.}
\sq{A loan of Rs 25000 is to be paid back in 6 bi-annual installments the first one starting after six month from the date of the loan. If the interest is charged at the rate of 18\% p.a compounded bi- annually on the unpaid principal. What will be amount of bi- annual installment?}
\sq{Mr. Saleem purchased a colour T.V set on installments each installment of Rs. 500 p.m, the first installment is paid at the time of purchase. If the interest charged is 24\% p.a compounded monthly. Find the cash price of the T.V if deal had to be completed in one year.}
\sq{New age cable, a cable manufacturing company is expected to pay Rs. 5.40 after the every six month on a share of the company. What is the present value of a share if money is worth 12\% p.a. compounded semi- annually?}
\sq{Define an equation.}
\sq{Define linear equation.}
\sq{Solve 9x=81.}
\sq{Solve 3x-3=6}
\sq{Solve 4-3=2x+7}
\sq{Solve the equation 9x+4=4x+29}
\sq{Solve 3x+9=12}
\sq{Solve for x if 2x+20-5x=x-6+9x}
\sq{Solve 2(x+5)-8(x-6)=10}
\sq{Solve $\dfrac{3x}{8}+5=17$}
\sq{Solve 3x+5=20}
\sq{Solve for 8x+3=5x+9}
\sq{Find x if x-7=3(x+5)}
\sq{Solve $\dfrac{1}{2x}+\dfrac{1}{4x}=4$}
\sq{Solve the equation $\dfrac{x}{3}-\dfrac{x}{6}=1$}
\sq{Find x if $\dfrac{x}{8}+\dfrac{x}{9}=72$}
\sq{Find the value of x if: 2x+7=9.}
\sq{Solve the equation 4x-3=2x+7}
\sq{Find the value of x if 3x+2=2x+6}
\sq{If 4x-3=2x+7 then find the value of x.}
\sq{Solve $\dfrac{7x+8}{3x+1}=\dfrac{5}{3}$}
\sq{Solve $\dfrac{3x+2}{4}=\dfrac{2x+6}{5}$}
\sq{Find x if $\dfrac{3x}{2}-\dfrac{x}{2}=\dfrac{5(x-4)}{6}$}
\sq{Find the value of x, $\dfrac{3x-1}{2-x}=2$}
\sq{Solve for x, $\dfrac{3}{x}+\dfrac{1}{5}=\dfrac{9}{2}$}
\sq{Solve the equation 2(x+5)-(x-6)=10}
\sq{Solve 3x-7=5x+2}
\sq{Solve 3x+5=7}
\sq{Solve $\dfrac{1}{x}+\dfrac{2}{x}=15$}
\sq{Solve 2(x-5)+8(x-3)=40}
\sq{Find x if 5(x-7)=3(x+5)-20}
\sq{Find x if 2x-7=13}
\sq{Find the value of x when 2x-4=8}
\sq{Find x if $\dfrac{3x}{8}-\dfrac{x}{2}=5$}
\sq{Solve form $x$. $\dfrac{1}{40}+\dfrac{1}{60}=\dfrac{1}{x}$}
\sq{Solve for x, 3(4x-2)=4(2x+3)}
\sq{Solve the equation $\dfrac{2x}{7}+2=0$}
\sq{Solve the equation $\dfrac{9}{x+4}=\dfrac{5}{x-8}$}
\sq{Solve for x, 5(x-3)=2(2x+4)}
\sq{Solve for x, 4x-7=2x+1}
\sq{Solve 2y+9=3}
\sq{Solve for x in the equation. 3(x-6)=2(x+3)+4}
\sq{Solve the equation. 2(4x+3)=3(2x+3)}
\sq{If $\dfrac{1}{4}$ of an amount of Rs. 60 what is the amount?}
\sq{If $\dfrac{1}{5}$ of an amount is Rs.10000 then find the amount.}
\sq{Nine times of numbe is 270. What is number.}
\sq{Solve the equation $\dfrac{1}{x}+\dfrac{2}{x}=15$}
\sq{Six times a number is 240. what is the number.}
\sq{When number is added to two third of itself the result is 45. Find the number.}
\sq{Eight times of a number is 480. what is the number?}
\sq{If seven times a number if 49. Find the number.}
\sq{Five times a number is 150. What is the number.}
\sq{Find two consecutive integers whose sum is 101.}
\sq{Find two consecutive integers whose sum is 41.}
\sq{Find two consecutive odd integers whose sum is 16.}
\sq{Find two consecutive odd integers whose sum is 24.}
\sq{Sum of two odd integers is 36.. find integers.}
\sq{Find three consecutive integers whose sum is 84.}
\sq{Find two consecutive integers whose sum is 65.}
\sq{Find two consecutive odd integers whose sum is 12.}
\sq{Find three consecutive even number whose sum is 228.}
\sq{A number with the addition of itself becomes 38.What is the number.}
\sq{Solve for x. 2x+20-5x=x-6+9x.}
\sq{If 9 times a number is 180. find the number.}
\sq{Solve $\dfrac{2x}{7}+1=0$}
\sq{Find two consecutive odd integers whose sum is 36.}
\sq{If x-4=16 , what does x equal?}
\sq{If y+16=20, what does y equal?}
\sq{If $\dfrac{1}{4}$ of an account if Rs. 150 find the amount.}
\sq{Six times a number is 180. what is the number?}
\sq{If 4x-6=2x+8 what does x equal?}
\sq{Solve $\dfrac{3x-4}{3}+\dfrac{x-3}{2}=\dfrac{7}{6}$}
\sq{Solve: 4(3y-9)=7(2-5y)+22y}
\sq{Find two consecutive odd integers, whose sum is 36.}
\sq{If a rectangle has a length that is 2 feet less then 3 times its width and its perimeter is 52 feets. Find its length \& width.}
\sq{How many gallons of 45\% acid solution and how many gallons of a 15\% acid solution shall be mixed to obtain 60 gallons of a 25\% acid solution?}
\sq{Define a quadratic equation.}
\sq{What is standard form of quadratic formula?}
\sq{Define quadratic formula.}
\sq{Write any two methods of finding solution of quadratic equation.}
\sq{What is discriminant?}
\sq{Express the following equation into standard for: $\dfrac{1}{x+3}-\dfrac{1}{x-3}=3$}
\sq{Express the following equation into standard for: $\dfrac{1}{x+4}-\dfrac{1}{x-4}=4$}
\sq{Express the following equation into standard for: $3x^{2}-1=\dfrac{11}{12x}$}
\sq{Solve the quadratic equation by factorization: $x^{2}-8x+15=0$}
\sq{Solve the quadratic equation by factorization: $2m^{2}-xm=0$}
\sq{Solve the quadratic equation by factorization: $x^{2}+5x+6=0$}
\sq{Solve the quadratic equation by factorization: $2x^{2}-9x+4=0$}
\sq{Solve the quadratic equation by factorization: $x^{2}+x-2=0$}
\sq{Solve the quadratic equation by factorization: $x^{2}-4x-21=0$}
\sq{Solve the quadratic equation by factorization: $42+Z-Z^{2}=0$}
\sq{Solve the quadratic equation by factorization: $x^{2}-5x+6=0$}
\sq{Solve the quadratic equation by factorization: $3x^{2}-10x+3=0$}
\sq{Solve the quadratic equation by factorization: $9x^{2}-12x-5=0$}
\sq{Solve the quadratic equation by factorization: $x^{2}+x-6=0$}
\sq{Solve the quadratic equation by factorization: $x^{2}-7x+10=0$}
\sq{Solve the quadratic equation by factorization: $2x^{2}+16=12x$}
\sq{Solve the quadratic equation by factorization: $x^{2}-16=0$}
\sq{Solve the quadratic equation by factorization: $x^{2}-11x+28=0$}
\sq{Solve the quadratic equation by factorization: $x^{2}-4x-21=0$}
\sq{Solve the quadratic equation by factorization: $x^{2}-x-2=0$}
\sq{Solve the quadratic equation by factorization: $x^{2}-5x+6=0$}
\sq{Solve the quadratic equation by factorization: $x^{2}-7x+12=0$}
\sq{Solve the quadratic equation by factorization: $x^{2}-4x+3=0$}
\sq{Solve the quadratic equation by factorization: $2x^{2}+7x+6=0$}
\sq{Solve the quadratic equation by factorization: $x^{2}-9x+14=0$}
\sq{Solve the quadratic equation by factorization: $x^{2}-8x+15=0$}
\sq{Solve the quadratic equation by factorization: $4x^{2}+5x-21=0$}
\sq{Solve the quadratic equation by factorization: $x^{2}+3x+2=0$}
\sq{Solve the following equation by completing square method: $x^{2}+5x-11=0$}
\sq{Solve the following equation by completing square method: $3x^{2}-2x=12$}
\sq{Solve the following equation by quadratic formula: $\dfrac{3}{x-2}+\dfrac{1}{x+2}=5$}
\sq{Solve the following equation by quadratic formula: $\dfrac{1}{x+3}-\dfrac{1}{x-3}=3$}
\sq{Solve the following equation by quadratic formula: $2x^{2}-x-8=0$}
\sq{Solve the following equation by quadratic formula: $3x^{2}+4x-5=0$}
\sq{Solve the following equation by quadratic formula: $x+\dfrac{1}{x}=4$}
\sq{Solve the following equation by quadratic formula: $3x^{2}-9x+5=0$}
\sq{Solve the following equation by quadratic formula: $4x^{2}+7x-1=0$}
\sq{Solve the following equation by quadratic formula: $2x^{2}-10x+5=0$}
\sq{Solve the following equation by quadratic formula: $x^{2}+x-3=0$}
\sq{Solve the following equation by quadratic formula: $2x^{2}-3x-6=0$}
\sq{Solve the following equation by quadratic formula: $x+\dfrac{1}{x}=1$}
\sq{Solve the following equation by quadratic formula: $x-\dfrac{2}{x}=1$}
\sq{Solve the following equation by quadratic formula: $x,2x^{2}+25x+72=0$}
\sq{Solve the following equation by quadratic formula: $y^{2}+40y+375=0$}
\sq{Solve the following equation by quadratic formula: $x^{2}-3x+2=0$}
\sq{Solve the following equation by quadratic formula: $9x^{2}-81=0$}
\sq{Find the discriminant of the following quardratic equation: $3x^{2}-5x+2=0$}
\sq{Find the discriminant of the following quadratic equation: $x^{2}+5x+2=0$}
\sq{Find the discriminant of the following quadratic equation $4x^{2}-3x+3=0$}
\sq{Find the discriminant of the following quadratic equation: $8x^{2}=72$}
\sq{Find the discriminant of the following quadratic equation: $x^{2}+10x+20=0$}
\sq{Find the discriminant of the following quadratic equation: $2x^{2}-3x-6=0$}
\sq{Find the discriminant of the following quadratic equation: $4x^{2}+13x+3=0$}
\sq{Find the discriminant of the following quadratic equation: $x^{2}-8x+15=0$}
\sq{Write down the sum roots of quadratic equation: $5x^{2}-6x+1=0$}
\sq{olve the equation $y^{2}+24=-10y$}
\sq{Solve the equation $x^{2}-8x+15=0$}
\sq{Find ther solution set of the equation $42+z-z^{2}=0$.}
\sq{Solve the equation $ax^{2}+bx+c=0$by completing square method.}
\sq{Solve the equation $\dfrac{x-1}{x}+\dfrac{x}{x+1}=\dfrac{6}{5}$}
\sq{Define reciprocal equation.}
\sq{Define an exponential equation.}
\sq{Define extraneous roots.}
\sq{Solve the equation for x if $4^{x}=16$.}
\sq{Solve $(5)^{1+x}+(5)^{1-x}=10$.}
\sq{Apply componendo and dividend rule on $\dfrac{\sqrt{x+1}+\sqrt{x-1}}{\sqrt{x+1}-\sqrt{x-1}}=\dfrac{2}{1}$}
\sq{Solve: $\dfrac{\sqrt{x-3}-\sqrt{x+3}}{\sqrt{x-3}+\sqrt{x+3}}=\dfrac{7}{4}$}
\sq{Solve: $\dfrac{1}{x+1}+\dfrac{2}{x+2}=\dfrac{4}{x+4}$}
\sq{Solve: $\sqrt{\left(\dfrac{x}{1-x}\right)}+\sqrt{\left(\dfrac{1-x}{x}\right)}=2\dfrac{1}{6}$}
\sq{Solve: $x^{\frac{2}{3}}-2x^{\frac{1}{3}}=0$}
\sq{Solve: $\sqrt{x+3}+\sqrt{x}=\sqrt{x+8}$}
\sq{Solve: $2x^{2}-4x-\sqrt{x^{2}-2x+13}=2$}
\sq{Solve: $(x-2)(x-4)(x-6)(x-8)=9$}
\sq{Solve: $x^{4}+4x^{3}+5x^{2}+4x+1=0$}
\sq{Solve: $\dfrac{\sqrt{x+2}+\sqrt{x-2}}{\sqrt{x+2}-\sqrt{x-2}}=\dfrac{5}{2}$}
\sq{Define simultaneous equation.}
\sq{Show that \{(30,8)\} is solution set of the system 3x+2y=106, 2x+4y=92.}
\sq{Solve 3y-4y=7 for y=2.}
\sq{Find the solution set of simultaneous equation x=y and 2x+y=3.}
\sq{Solve the following system of simultaneous equation. 2x+y=8, x-y=1}
\sq{Solve the following system of simultaneous equation. 4x-2y=y and 5x- y=10.}
\sq{Solve the following system of simultaneous equation. 2x-y=50 , x+y=4}
\sq{Solve the following system of simultaneous equation. x+y=40, x-y=15}
\sq{Solve the following system of simultaneous equation. x+y=16, x-y=2}
\sq{Solve the following system of simultaneous equation. 5x+7y=11, 2x- 4y=-16}
\sq{Solve the following system of simultaneous equation. 5x+2y=64 and 2x-y=4.}
\sq{Solve the following system of simultaneous equation. x+y=12 , x-y= 8}
\sq{Solve the following system of simultaneous equation. 2x+7y=12, 2x+y=-6}
\sq{Solve the following system of simultaneous equation. 4y-x=5, y+x=5}
\sq{Solve the following system of simultaneous equation. x-y=6 and 2x+y=3}
\sq{Solve the following system of simultaneous equation. x+4y=14 and 3x-2y=7}
\sq{Solve the following system of simultaneous equation. x,y if 2(x- 4)=3(y-3) and y-9x=13}
\sq{Solve the following system of simultaneous equation. x and y 4y- 2x=50 and 5x-y=10.}
\sq{Solve the following system of simultaneous equation. x and y x- 3y=-6, 3x-2y=-4}
\sq{Solve the following system of simultaneous equation. x+y=8, x-y=4}
\sq{Solve the following system of simultaneous equation. x+y=8, x-y=4}
\sq{Solve the following system of simultaneous equation. 3x+2y=10 ; 2x-2y=5}
\sq{Solve the following system of simultaneous equation. x-2y=1; 3x+2y=5}
\sq{Solve the following system of simultaneous equation. x-2y=1; 3x+2y=7}
\sq{Solve the following system of simultaneous equation. x+y=3, 2x-y=6}
\sq{Solve the following system of simultaneous equation. 5x+2y=33 and 5x-2y=4}
\sq{Solve the following system of simultaneous equation. x-y=2 and x+y=16}
\sq{Solve the following system of simultaneous equation. 4x-y=5, x+y=5}
\sq{Find two positive numbers having a sum 23 and a product 126.}
\sq{The sum of number os 42 and their difference is 14. Find the number.}
\sq{Find x and y from the following set of equations: 4y-2x=50 and 5x- y=10}
\sq{Solve 3x+2y=54 and 2x-3y=10}
\sq{Find two numbers such that $\dfrac{1}{2}$ of the first point plus $\dfrac{1}{3}$ of the second is equal to 45, and $\dfrac{1}{2}$ of the second plus $\dfrac{1}{5}$ of the first is equal to 40.}
\sq{A man has invested Rs.220000 and on it he earned Rs. 12,200 in last year. Part of the money is invested at 5\% interest and the other part at 7\%. How much he invested in each part?}
\sq{Define a functions.}
\sq{Decide whether the eqaution $x^{2}+y^{2}=1$ is a function.}
\sq{Define domain of the function.}
\sq{Define range of a function.}
\sq{Define independent variable.}
\sq{Define implicit function with example.}
\sq{Define explicit function.}
\sq{Define quadratic function.}
\sq{Write polynomial function of degree one.}
\sq{Define even function and give and example.}
\sq{Define odd function.}
\sq{Define variable and constant.}
\sq{Define constant function.}
\sq{Define domain and range the following function/ relation. \{(1,3), (1,-1),(1,4),(1,-2)\}.}
\sq{Define domain and range the following function/ relation. \{(0,0),)1,2), (1,3),(2,4)\}.}
\sq{Define domain and range the following function/ relation. \{(1,3),(3,3), (5,1),(6,1)\}}
\sq{Define domain and range the following function/ relation. \{(-2,-3), (-1,1),(0,1)(1,3)(2,5)\}}
\sq{Define domain and range the following function/ relation. \{(2,3), (3,7),(4,8),(5,11)\}.}
\sq{Define domain and range the following function/ relation. $f(x)=\sqrt{2x-3}$}
\sq{Define domain and range the following function/ relation. $f(x)=\dfrac{x}{5}$ find domain}
\sq{Define domain and range the following function/ relation. f(x)=x+1.}
\sq{Define domain and range the following function/ relation. $f(x)=\dfrac{1}{x+4}$. Find domain?}
\sq{Define domain and range the following function/ relation. $f(y)=\dfrac{1}{y-1}$. Find domain?}
\sq{Define domain and range the following function/ relation. $f(x)=\dfrac{7}{x-1}$. Find domain?}
\sq{Define domain and range the following function/ relation. $f(x)=\dfrac{1}{x-3}$. Find domain?}
\sq{Define domain and range the following function/ relation. $f(x)=\sqrt{x-3}$}
\sq{Define domain and range the following function/ relation. Find domain of: $f(x)=x+\dfrac{1}{x}$}
\sq{Determine whether given function f(x) is (a) even (b) odd (c) neither even nor odd. f(x)= $x^{2}$-x+5}
\sq{Determine whether given function f(x) is (a) even (b) odd (c) neither even nor odd. f(x)= $x^{3}$-x is and odd function.}
\sq{Determine whether given function f(x) is (a) even (b) odd (c) neither even nor odd. f(x)= $x^{5}$+$x^{3}$ is and odd function.}
\sq{Determine whether given function f(x) is (a) even (b) odd (c) neither even nor odd. f(x)= $3x^{3}$-x is an even / odd function.}
\sq{Determine whether given function f(x) is (a) even (b) odd (c) neither even nor odd. f(x)= $x^{3}$+2x+1}
\sq{Determine whether given function f(x) is (a) even (b) odd (c) neither even nor odd. f(x)= $x^{3}$+5x is an odd function.}
\sq{Determine whether given function f(x) is (a) even (b) odd (c) neither even nor odd. f(x)= $x^{2}$+5 is even or odd.}
\sq{Determine whether given function f(x) is (a) even (b) odd (c) neither even nor odd. f(x)= $x^{4}$+$x^{2}$ is an even or odd function.}
\sq{Determine whether given function f(x) is (a) even (b) odd (c) neither even nor odd. f(x)= $x^{3}$-3 is netither even nor odd function.}
\sq{If f(x)=10x+5 then find the values of f(1) and f(-1).}
\sq{If f(x) = $x^{2}$+3, then find f(2), f(3).}
\sq{If $f(x)=\dfrac{x}{x+1}$ then find $f(-2)$}
\sq{If Z=$X^{2}$-$y^{2}$, find the Z when X=1, Y=5.}
\sq{If f(x)=$x^{2}$-4 then find the value of f(4) and $f(\sqrt{2})$}
\sq{If y=3x+4 find y for x =0, x=1.}
\sq{If y=40-4x, find x when y=-4.}
\sq{If f(x) = 2x+1, find f(-4),f(3).}
\sq{If the function g(u)=$u^{2}$+u then find g(2v) and g(-$x^{2}$)}
\sq{If f(x) = $x^{4}$-$x^{2}$+3 then find $f(-\sqrt{2})$}
\sq{If f(x) = $x^{2}$-2x+3 then find x=0 and x=2.}
\sq{If f(x)= 2x+3 find f(0) and f(1)}
\sq{If $f(x)=\dfrac{x}{x+1}$, find f(-2) and f(1).}
\sq{If f(x)=2$x^{2}$-x+4 find f(-3) and $f\left(\dfrac{1}{2}\right)$}
\sq{If f(x)= -2x+4, then find $f\left(-\dfrac{1}{2}\right)$ and $f\left(\dfrac{5}{2}\right)$}
\sq{If f(x)=4-2x. find f(-3).}
\sq{Find the quantity demand for the following function Q=5+2P when price P-5,10.}
\sq{If H(S)-=$S^{2}$-4, find H(3) and $H(\sqrt{2})$}
\sq{If f(x) = $x^{2}$-2x-3 then find f(-1) and f(0).}
\sq{If f(x)= $x^{2}$-2x+0 then find f(-1)}
\sq{if $f(x)=\dfrac{x^{2}+3x-2}{x+4}$ then find f(0).}
\sq{If f(x)=5+2x-$x^{2}$, find f(-3),f(1).}
\sq{Find f(x) if f(x-1)=3$x^{2}$+2x-5.}
\sq{If Z=$x^{2}$-Y then find Z, when z=4, y=3.}
\sq{If f(x)= $x^{2}$-2x+1 then find f(a+b)}
\sq{if $f(x)=\sqrt{3x-4}$ find f(2) and f(0).}
\sq{If f(x)= $x^{2}$-2x+1, then find f(-1) and f(-2).}
\sq{Find the domain of the following function: $g(x)=\dfrac{1}{x-5}$}
\sq{Find the domain of the following function: $h(z)=\sqrt{2z-3}$}
\sq{Find the domain of the following function: $\phi(t)=\dfrac{1}{2}t-4$}
\sq{If $g(x)=\dfrac{1}{x-5}$, then find the value of g(x) for x=-2,-1,0 and 6.}
\sq{Find Find H(4),$H(\sqrt{2})$ and $H\left(\dfrac{2}{3}\right)$ from a fraction $H(s)=s^{2}-3$ [text cut off at the page edge in the source]}
\sq{If given function is g(u)=$u^{2}$+u. (i) g(-2)=? (ii) g(2v)=? and (iii) g(-$x^{2}$)=?}
\sq{Develop linear function/equation from orderes paris (x,y):(2,8) and (5,23).}
\sq{Develop quadratic function $y=ax^{2}+bx+c$ using following set of information:
\par\smallskip\hspace*{11em}\begin{tabular}{|l@{\hspace{3em}}c@{\hspace{3em}}c@{\hspace{3em}}c|}\hline \textbf{x} & -3 & -1 & 2 \\ \textbf{y} & 0 & -4 & 5 \\ \hline\end{tabular}}
\sq{What is rectangular coordinate system?}
\sq{Write two uses of coordinate system.}
\sq{What is the equation of x-axis?}
\sq{What is the equation of y-axis.}
\sq{What are the coordinate axes?}
\sq{What is the graph of quadratic function?}
\sq{What is the graph of linear function?}
\sq{Define origin.}
\sq{Define x-intercept.}
\sq{Define y-intercept.}
\sq{What is x-intercept of a function?}
\sq{What is y-intercept of a function?}
\sq{Define Break Even point.}
\sq{In how many regions a plane is divided by coordinate axes? what we call there regions?}
\sq{Identify the location of following points on a x,y- plane: (a) (3,2), (b) (1/2,-3), (c) (-2,4) (d) (-3,-2) (e) (0,0) (f) (0,1/2) (g) (10,-7).}
\sq{Find x-intercept of the function 2x+3y=5}
\sq{If z=$x^{2}$-$y^{2}$ , find z when z=1, y=5.}
\sq{Find x and y intercepts of 3y+x=12.}
\sq{If y=10-7x , find the value of y for x=1.1.}
\sq{What is absissa of the point on the graph of: 3x+4y+2=0 if ordinate is 4?}
\sq{Find the x-intercept and y-intercept of the equation 12y-15x-14=0.}
\sq{Find x-intercept of the line given by equation y-5x+3=0.}
\sq{Find y-intercept of 2x+y=8.}
\sq{If y=2x+1 find slop and y-intercept.}
\sq{Find the slop and y-intercept in the equation y-5x+3=0.}
\sq{Find the intercept of the line 2x+y=5.}
\sq{Find an equation of straight line whose slop is -2 and y-intercept is $\dfrac{4}{5}$.}
\sq{Find slop and y intercept of the line y-5x+7=0.}
\sq{Identify the intercept of the graph of the equation 2x-3y=6.}
\sq{Conver 2x-3y=6 into intercepts.}
\sq{Draw the graph of f(x) = $x^{2}$}
\sq{Draw graph of linear function 3x=-y+2.}
\sq{Draw graph of f(x)=x.}
\sq{Find the x-intercepts and y-intercepts of the graph of y=2x+4 and sketch the graph.}
\sq{Swift pens has daily fixed cost of Rs. 3000 as a result of salaries and building operations. Each ballpoint produced costs Rs. 5 and is sold of Rs. 7 what is the beak-even point? That is , how many ballpoints must be sold each day to get no loss and no profit?}
\sq{A producer sells items for Rs.0.30 each . If his cost of production is C= 0.15 X + 105 where X is the number of its product sold. Find his break even point if he can change his cost. C= 0.12 X + 110 would it be advantageous to him?}
\sq{Define matrices.}
\sq{Define order of matrix.}
\sq{Define row matrix.}
\sq{Define a column matrix. Give example.}
\sq{Define regular matrix.}
\sq{Define square matrix.}
\sq{Define null matrix.}
\sq{Define diagonal matrix.}
\sq{Define scalar matrix.}
\sq{Define unit or multiplicative identity of matrix.}
\sq{Name any three types of matrices.}
\sq{If A = $\begin{bmatrix}1&4\\3&2\\2&5\end{bmatrix}$ and B = $\begin{bmatrix}6&5\\2&6\\1&1\end{bmatrix}$ find A+B.}
\sq{If A = $\begin{bmatrix}2&3\\1&4\end{bmatrix}$, B$\begin{bmatrix}1&3\\4&-3\end{bmatrix}$ find A-B.}
\sq{If A = $\begin{bmatrix}0&2\\3&5\end{bmatrix}$, B = $\begin{bmatrix}4&1\\2&3\end{bmatrix}$. Find A+B and A-B.}
\sq{If A = $\begin{bmatrix}2&3\\-2&1\end{bmatrix}$, B = $\begin{bmatrix}0&-4\\1&5\end{bmatrix}$. Find A+B.}
\sq{Find A+B if A = $\begin{bmatrix}2&-1&3\\-5&4&7\end{bmatrix}$ B = $\begin{bmatrix}1&5&7\\9&3&10\end{bmatrix}$.}
\sq{Find A+B where A = $\begin{bmatrix}1&4\\3&2\\2&5\end{bmatrix}$ B = $\begin{bmatrix}6&5\\2&6\\1&1\end{bmatrix}$}
\sq{[question text missing in the source]}
\sq{If A = $\begin{bmatrix}10&3&5\\8&11&2\\9&6&7\end{bmatrix}$, B = $\begin{bmatrix}9&7&3\\2&5&8\\4&12&2\end{bmatrix}$ find A+B.}
\sq{If A = $\begin{bmatrix}6&-3\\-4&5\end{bmatrix}$ and B = $\begin{bmatrix}2&3\\1&4\end{bmatrix}$ then find A-B.}
\sq{If A = $\begin{bmatrix}2&1\\1&1\end{bmatrix}$ and B = $\begin{bmatrix}1&1\\0&1\end{bmatrix}$ the find 3A-2B.}
\sq{If A = $\begin{bmatrix}1&2\\-1&2\end{bmatrix}$, B = $\begin{bmatrix}0&2\\1&3\end{bmatrix}$ find 2A+3B.}
\sq{If A = $\begin{bmatrix}6\\7\end{bmatrix}$ and B = $\begin{bmatrix}1\\0\end{bmatrix}$ find A-3B.}
\sq{If A = $\begin{bmatrix}1&2\\3&4\end{bmatrix}$, B = $\begin{bmatrix}4&3\\5&2\end{bmatrix}$ find 2A+3B.}
\sq{If A = $\begin{bmatrix}4&7\\4&15\end{bmatrix}$ and B = $\begin{bmatrix}3&8\\2&20\end{bmatrix}$ Compute (a) 2A-B (b) A+2B}
\sq{If C = $\begin{bmatrix}0&8\end{bmatrix}$ find $\dfrac{1}{2}$C.}
\sq{Find AB if A = $\begin{bmatrix}3&4\end{bmatrix}$ and B = $\begin{bmatrix}4\\5\end{bmatrix}$}
\sq{Find the value of X if $\begin{bmatrix}3&5\\-5&4\end{bmatrix}$+X = $\begin{bmatrix}1&0\\0&1\end{bmatrix}$}
\sq{Find AB if A = $\begin{bmatrix}1&5\end{bmatrix}$, B = $\begin{bmatrix}1\\5\end{bmatrix}$. BA.}
\sq{Let A = $\begin{bmatrix}5\\7\\3\end{bmatrix}$ and B = $\begin{bmatrix}2&-1&0\end{bmatrix}$ If possible find at AB and BA}
\sq{If A = $\begin{bmatrix}1&2&3\end{bmatrix}$ and B = $\begin{bmatrix}1\\2\\3\end{bmatrix}$ then find AB and BA.}
\sq{If A = $\begin{bmatrix}1&2\\0&3\end{bmatrix}$, B$\begin{bmatrix}0&2\\1&4\end{bmatrix}$ Find AB.}
\sq{If A = $\begin{bmatrix}1&0\\0&1\end{bmatrix}$, find $A^{2}$}
\sq{If A = $\begin{bmatrix}1&2\\3&4\end{bmatrix}$, B = $\begin{bmatrix}6&8\\4&-4\end{bmatrix}$ then find AB?}
\sq{If A = $\begin{bmatrix}1&2\\-1&3\end{bmatrix}$, B = $\begin{bmatrix}-2\\5\end{bmatrix}$ then find AB.}
\sq{If P = $\begin{bmatrix}2&1\\3&1\end{bmatrix}$, Q = $\begin{bmatrix}1&0\\0&1\end{bmatrix}$ then find PQ.}
\sq{If A = $\begin{bmatrix}3&1\\2&2\end{bmatrix}$, find $A^{2}$}
\sq{If A = $\begin{bmatrix}4&3\\2&5\end{bmatrix}$, find $A^{2}$}
\sq{If A = $\begin{bmatrix}0&1&-1\\4&-3&4\\3&-3&4\end{bmatrix}$ find $A^{2}$.}
\sq{If A = $\begin{bmatrix}4&5\\2&3\end{bmatrix}$ Find $A^{2}$.}
\sq{If A = $\begin{bmatrix}3&1\\2&0\end{bmatrix}$ and $A^{2}$=AB then find matrix B.}
\sq{If A = $\begin{bmatrix}2&3\\1&6\end{bmatrix}$ and B = $\begin{bmatrix}1&7\\6&3\end{bmatrix}$ Find AB.}
\sq{If A = $\begin{bmatrix}1&-2\\3&-4\end{bmatrix}$ and B = $\begin{bmatrix}8&7\\6&5\end{bmatrix}$ Find AB.}
\sq{Define transpose of a matrix.}
\sq{Define symmetric matrix.}
\sq{Define skew symmetric matrix.}
\sq{If A = $\begin{bmatrix}4&1\\3&2\end{bmatrix}$ find $A^{t}$}
\sq{If A = $\begin{bmatrix}4&5\\2&3\end{bmatrix}$ find $A^{t}$}
\sq{Find transpose of A = $\begin{bmatrix}3&2&4\\2&3&7\end{bmatrix}$}
\sq{Find A if 2A + $\begin{bmatrix}1&2\\4&6\end{bmatrix}$ = 0}
\sq{Find the value of x from x + $\begin{bmatrix}3\\4\end{bmatrix}$ = $\begin{bmatrix}5\\2\end{bmatrix}$}
\sq{Find matrix x if 3x + $\begin{bmatrix}1&2\\3&4\end{bmatrix}$ = $\begin{bmatrix}-1&-2\\1&2\end{bmatrix}$}
\sq{Find A if 3A - $\begin{bmatrix}6&-9&15\\-3&18&0\end{bmatrix}$}
\sq{Find B if 2B + $\begin{bmatrix}2&5\\4&6\end{bmatrix}$ = 0}
\sq{A = $\begin{bmatrix}1&-1&2\\2&1&0\end{bmatrix}$, B = $\begin{bmatrix}1&2\\2&0\end{bmatrix}$ then find AB.}
\sq{If A = $\begin{bmatrix}3&1\\2&0\end{bmatrix}$, B = $\begin{bmatrix}4&-1\\2&3\end{bmatrix}$ find BA.}
\sq{If A = $\begin{bmatrix}3&8\\5&4\end{bmatrix}$, find $A^{2}$.}
\sq{Find the product BA, if A = $\begin{bmatrix}3&-2\\1&0\end{bmatrix}$ and B = $\begin{bmatrix}1&2\\3&1\end{bmatrix}$}
\sq{If A = $\begin{bmatrix}2&3\\1&5\end{bmatrix}$, B = $\begin{bmatrix}2&3&1\\1&1&1\end{bmatrix}$ then find AB.}
\sq{A = $\begin{bmatrix}1\\5\\4\end{bmatrix}$ and B = $\begin{bmatrix}1&8&7\end{bmatrix}$, find AB.}
\sq{Find the value of x and y if $\begin{bmatrix}x+3&1\\-3&3y-4\end{bmatrix}$ $\begin{bmatrix}y&1\\-3&2x\end{bmatrix}$}
\sq{Show that A = $\begin{bmatrix}0&a\\-a&0\end{bmatrix}$ is skew-symmetric matrix.}
\sq{Find x and y if $\begin{bmatrix}x-2&5\\7&y-7\end{bmatrix}$ = $\begin{bmatrix}6&5\\7&2\end{bmatrix}$}
\sq{Define lower triangular matrix.}
\sq{If A = $\begin{bmatrix}1&2\\3&4\end{bmatrix}$ and B = $\begin{bmatrix}6&8\\4&-4\end{bmatrix}$, Find A+B}
\sq{If A = $\begin{bmatrix}1&4&5\\6&7&8\end{bmatrix}$, find 2A}
\sq{If A = $\begin{bmatrix}1&2\\3&4\end{bmatrix}$ and B = $\begin{bmatrix}6&8\\4&-4\end{bmatrix}$ Find A-B}
\sq{If A = $\begin{bmatrix}2&3\\1&6\end{bmatrix}$ and B = $\begin{bmatrix}1&7\\6&3\end{bmatrix}$, then find AB, BA}
\sq{If A = $\begin{bmatrix}2&3\\1&5\end{bmatrix}$ and C = $\begin{bmatrix}2&3&1\\1&1&1\end{bmatrix}$, Find AC and CA if possible.}
\sq{If A = $\begin{bmatrix}4&5\\-2&3\end{bmatrix}$, find $A^{2}$}
\sq{Show that $\begin{bmatrix}-1&2&3\\2&1&0\\3&5&-1\end{bmatrix}$$\begin{bmatrix}a\\b\\c\end{bmatrix}$ = $\begin{bmatrix}-a+2b+3c\\2a+b\\3a+5b-c\end{bmatrix}$}
\sq{If A = $\begin{bmatrix}1&-1&2\\2&1&0\end{bmatrix}$ and B = $\begin{bmatrix}1&2\\2&3\\-1&1\end{bmatrix}$, Then verify that $(AB)^{t}=B^{t}A^{t}$}
\sq{If A is a square matrix, then show that A+A$^{t}$ is symmetric?}
\sq{Define the determinant?}
\sq{What do you understand by the singular matrix.}
\sq{Define non-singular matrix.}
\sq{Define adjoint of the matrix.}
\sq{Define inverse matrix.}
\sq{Find the determine of A is: A = $\begin{bmatrix}a&b\\c&d\end{bmatrix}$}
\sq{Find the determine of A is: A = $\begin{bmatrix}4&-4\\6&-5\end{bmatrix}$}
\sq{Find the determine of A is: A = $\begin{bmatrix}-3&-2\\4&5\end{bmatrix}$}
\sq{Find the determine of A is: A = $\begin{bmatrix}2&9\\3&8\end{bmatrix}$}
\sq{Find the determine of A is: A = $\begin{bmatrix}1&2\\0&3\end{bmatrix}$}
\sq{Find the determine of A is: A = $\begin{bmatrix}2&0&0\\2&4&1\\5&6&7\end{bmatrix}$}
\sq{Find the determine of A is: A = $\begin{bmatrix}2&-4\\1&5\end{bmatrix}$}
\sq{If A = $\begin{bmatrix}3&4\\5&6\end{bmatrix}$, find $\dfrac{1}{2}|A|$}
\sq{Show that following matrix is singular: A = $\begin{bmatrix}1&2\\2&4\end{bmatrix}$}
\sq{Show that following matrix is singular: A = $\begin{bmatrix}2&-3\\2&-3\end{bmatrix}$}
\sq{Find x if the following matrix is singular: A = $\begin{bmatrix}6&3\\4&x\end{bmatrix}$}
\sq{Find x if the following matrix is singular: A = $\begin{bmatrix}3&x\\-5&10\end{bmatrix}$}
\sq{Find x if the following matrix is singular: A = $\begin{bmatrix}1&-2\\-3&x\end{bmatrix}$}
\sq{Find x if the following matrix is singular: A = $\begin{bmatrix}3&x\\4&8\end{bmatrix}$}
\sq{Find x if the following matrix is singular: A = $\begin{bmatrix}2&x\\5&10\end{bmatrix}$}
\sq{Find value of $\lambda$ if A = $\begin{bmatrix}\lambda&4\\3&2\end{bmatrix}$ si singular/}
\sq{Find x if $\begin{vmatrix}x&4\\3&2\end{vmatrix}$ = 0}
\sq{Find x if $\begin{vmatrix}x+3&x\\x+2&x-3\end{vmatrix}$ = 0}
\sq{Find adjoint of A of: A = $\begin{bmatrix}-1&-2\\3&4\end{bmatrix}$}
\sq{Find adjoint of A of: A = $\begin{bmatrix}2&1\\0&-5\end{bmatrix}$}
\sq{Find adjoint of A of: A = $\begin{bmatrix}1&2\\-1&3\end{bmatrix}$}
\sq{Find adjoint of A of: A = $\begin{bmatrix}1&0\\0&3\end{bmatrix}$}
\sq{Find adjoint of A of: A = $\begin{bmatrix}3&-3\\4&1\end{bmatrix}$, Find |A|.}
\sq{Find adjoint of A of: A = $\begin{bmatrix}6&2\\4&1\end{bmatrix}$, Find |A|.}
\sq{Find the inverse matrix (A$^{-1}$) if: [matrix missing in the source]}
\sq{Find the inverse matrix (A$^{-1}$) if: A = $\begin{bmatrix}1&2\\-1&2\end{bmatrix}$}
\sq{Find the inverse matrix (A$^{-1}$) if: A = $\begin{bmatrix}3&-1\\4&5\end{bmatrix}$}
\sq{Find the inverse matrix (A$^{-1}$) if: A = $\begin{bmatrix}2&3\\1&6\end{bmatrix}$}
\sq{Find the inverse matrix (A$^{-1}$) if: A = $\begin{bmatrix}-1&-2\\-3&6\end{bmatrix}$}
\sq{Find the inverse matrix (A$^{-1}$) if: A = $\begin{bmatrix}4&-7\\8&11\end{bmatrix}$}
\sq{Find the inverse matrix (A$^{-1}$) if: A = $\begin{bmatrix}4&6\\2&1\end{bmatrix}$}
\sq{Find the inverse matrix (A$^{-1}$) if: A = $\begin{bmatrix}1&3\\0&-1\end{bmatrix}$}
\sq{Find the inverse matrix (A$^{-1}$) if: A = $\begin{bmatrix}7&8\\4&5\end{bmatrix}$}
\sq{Find the inverse matrix (A$^{-1}$) if: A = $\begin{bmatrix}4&9\\7&6\end{bmatrix}$}
\sq{Find the inverse matrix (A$^{-1}$) if: A = $\begin{bmatrix}-2&6\\4&7\end{bmatrix}$}
\sq{Find the inverse matrix (A$^{-1}$) if: A = $\begin{bmatrix}3&2\\1&4\end{bmatrix}$}
\sq{Find the inverse matrix (A$^{-1}$) if: A = $\begin{bmatrix}4&1\\5&8\end{bmatrix}$}
\sq{Solve the following equation linear x+y=11 and x-y=7 by matrix method.}
\sq{Is the following matrix is singular or non-singular? $\begin{bmatrix}4&-4\\6&6\end{bmatrix}$}
\sq{Is the following matrix is singular or non-singular? $\begin{bmatrix}2&\sqrt{2}\\12&\sqrt{72}\end{bmatrix}$}
\sq{Find the value of x if $\begin{bmatrix}3&x\\-5&10\end{bmatrix}$ is a singular matrix.}
\sq{Find adjoint matrix of $\begin{bmatrix}0&-8\\4&\sqrt{2}\end{bmatrix}$.}
\sq{Find inverse of the following matrix: $\begin{bmatrix}4&9\\7&6\end{bmatrix}$}
\sq{Find inverse of the following matrix: $\begin{bmatrix}1&0\\0&1\end{bmatrix}$}
\sq{Evaluate the determinant $\begin{vmatrix}1&2&3\\4&5&6\\7&8&9\end{vmatrix}$}
\sq{Find the inverse of the following matrix: A = $\begin{bmatrix}2&1&2\\3&0&3\\-1&4&-5\end{bmatrix}$}
\sq{Solve: 2x+3y=10, 4x+8y=24.}
\sq{Solve the system of the equations: 2x-4y=16, 9x-18y=120}
\sq{Solve the following simultaneous equations by Cramer's rule: 3x+2y=12, x+5y=17}
\sq{What is decimal system.}
\sq{Define binary number system.}
\sq{Write digit on which decimal number system is based.}
\sq{Define number system used in computer.}
\sq{What are the meaning of 0 and 1 for computer?}
\sq{Convert the following decimal number into binary number system. 8}
\sq{Convert the following decimal number into binary number system. 12}
\sq{Convert the following decimal number into binary number system. 13}
\sq{Convert the following decimal number into binary number system. 14}
\sq{Convert the following decimal number into binary number system. 15}
\sq{Convert the following decimal number into binary number system. 19}
\sq{Convert the following decimal number into binary number system. 17}
\sq{Convert the following decimal number into binary number system. 21}
\sq{Convert the following decimal number into binary number system. 22}
\sq{Convert the following decimal number into binary number system. 23}
\sq{Convert the following decimal number into binary number system. 25}
\sq{Convert the following decimal number into binary number system. 26}
\sq{Convert the following decimal number into binary number system. 27}
\sq{Convert the following decimal number into binary number system. 28}
\sq{Convert the following decimal number into binary number system. 35}
\sq{Convert the following decimal number into binary number system. 49}
\sq{Convert the following decimal number into binary number system. 82}
\sq{Convert the following decimal number into binary number system. 90}
\sq{Convert the following decimal number into binary number system. 93}
\sq{Convert the following decimal number into binary number system. 99}
\sq{Convert the following decimal number into binary number system. 245}
\sq{Convert the following decimal number into binary number system. 247}
\sq{Convert the following decimal number into binary number system. 786}
\sq{Convert the following binary number into decimal number. (001)$_{2}$}
\sq{Convert the following binary number into decimal number. (101)$_{2}$}
\sq{Convert the following binary number into decimal number. (1001)$_{2}$}
\sq{Convert the following binary number into decimal number. (1011)$_{2}$}
\sq{Convert the following binary number into decimal number. (1101)$_{2}$}
\sq{Convert the following binary number into decimal number. (1110)$_{2}$}
\sq{Convert the following binary number into decimal number. (10001)$_{2}$}
\sq{Convert the following binary number into decimal number. (11010)$_{2}$}
\sq{Convert the following binary number into decimal number. (11101)$_{2}$}
\sq{Convert the following binary number into decimal number. (11111)$_{2}$}
\sq{Convert the following binary number into decimal number. $\dfrac{7}{8}$}
\sq{Convert (7)$_{10}$ to its equivalent number with binary system.}
\sq{Convert (99)$_{10}$ to its equivalent number with base 2.}
\sq{Convert 749 to its equivalent number with base 2.}
\sq{Convert (1100011)$_{2}$ to ots equivalent number with base 10.}
\sq{Change a binary number 1011101101 to its equivalent number with base 10.}
\sq{Convert $\dfrac{7}{8}$ to a binary number.}
\sq{Convert $\dfrac{39}{8}$ in binary system.}
\sq{Convert (101101.01)$_{2}$ into a decimal number.}
\sq{Convert (1010110.1111)$_{2}$ into a decimal number.}
\sq{convert 45.25 into a binary number.}
\sq{Convert 15.35 into a binary number.}
\sq{Evaluate (111)$_{2}$+(101)$_{2}$.}
\sq{Evaluate (1111)$_{2}$+(111)$_{2}$.}
\sq{Find the sum: (1101)$_{2}$+(1011)$_{2}$.}
\sq{Simplify: (111)$_{2}$+(100)$_{2}$.}
\sq{Find (11011)$_{2}$+(11111)$_{2}$.}
\sq{Find sum of (10101)$_{2}$ and (1101)$_{2}$.}
\sq{Find sum of: (11)$_{2}$ and (101)$_{2}$.}
\sq{Simplify: (1001)$_{2}$+(11)$_{2}$.}
\sq{Find the sum of: (1101)$_{2}$+(1100).}
\sq{Find (11111)$_{2}$+(10011)$_{2}$.}
\sq{Simplify: (1001)$_{2}$+(11)$_{2}$.}
\sq{Find the sum of: (1011)$_{2}$+(111)$_{2}$.}
\sq{Simplify (10110)$_{2}$+(1000)$_{2}$.}
\sq{Find (11101)$_{2}$+(111)$_{2}$.}
\sq{Add the binary numbers (1011)$_{2}$ and (1001)$_{2}$.}
\sq{Add the binary numbers (1011)$_{2}$ and (1001)$_{2}$.}
\sq{Add (10010110)$_{2}$ and (1011)$_{2}$}
\sq{Evaluate: (1110)$_{2}$-(101)$_{2}$.}
\sq{Subtract (1010)$_{2}$ from (110001)$_{2}$.}
\sq{Simplify: (1111)$_{2}$-(1010)$_{2}$.}
\sq{Subtract (111001)$_{2}$-(1001)$_{2}$.}
\sq{Evaluate: (10111)$_{2}$-(1111)$_{2}$.}
\sq{Find: (11101)$_{2}$-(111)$_{2}$.}
\sq{Subtract (11)$_{2}$ from (1101)$_{2}$.}
\sq{Subtract (1011100)$_{2}$ and (111100)$_{2}$}
\sq{olve: (10110)$_{2}$-(1000)$_{2}$.}
\sq{Simplify (1100)$_{2}$-(111)$_{2}$.}
\sq{Evaluate (11)$_{2}$ x (11)$_{2}$.}
\sq{Evaluate (101)$_{2}$ x (11)$_{2}$.}
\sq{Simplify: (1001)$_{2}$ x (101)$_{2}$.}
\sq{Simplify (1111)$_{2}$ x (1111)$_{2}$.}
\sq{Evaluate (111)$_{2}$ x (111$_{2}$).}
\sq{Find the product of (1011)$_{2}$ x (11)$_{2}$.}
\sq{Multiply the binary number (1011)$_{2}$ and (101)$_{2}$.}
\sq{Simplify: (1011)$_{2}$ x (111)$_{2}$.}
\sq{Evaluate: \{(100111)$_{2}$+(10101)$_{2}$\}-(10111)$_{2}$.}
\sq{Divide (100011)$_{2}$ by (101)$_{2}$.}
\sq{Compute: (1101)$_{2}$ x (10)$_{2}$.}
\sq{Find the sum (11)$_{2}$ and (101)$_{2}$.}
\sq{Add (10110)$_{2}$ and (1000)$_{2}$.}
\sq{Find the value of (10111011)$_{2}$+(101101)$_{2}$.}
\sq{Subtract 45 from 169 both in decimal as well as in binary system.}
\sq{Subtract 912 from 1942 in decimal and binary system.}
\sq{Compute the following: 1101 x 10}
\sq{Compute the following: 111 x 101}
\sq{Multiply 39 by 25 in both binary and decimal system.}
\sq{Perform the following binary multiplication and check the answers in the decimal system. 1101 x 111}
\sq{Perform the following binary multiplication and check the answers in the decimal system. 11011 x 10001}
\sq{Divide 87 by 14 in decimal and binary system by using repeated subtraction method.}
\sq{Divide the following number by subtracting and check by comparing in decimal system: $1101\div11$}
\sq{Divide the following number by subtracting and check by comparing in decimal system: $1101\div101$}
\sq{Simplify $\{(11101111)_{2}-(10001)_{2}\}-\{(1111000)_{2}-(1000$ [text cut off at the page edge in the source]}
\sq{Find value of the following by changing into decimal system. $\{(111100000)_{2}+(1110011111)_{2}\}-(29)_{10}$}
\sq{Find value of the following by changing into decimal system. $\{(1111111110)_{2}-(801)_{10}\}+\{(11100)_{2}-(1011)_{2}\}$}

\end{document}
